\chapter{雷达来源数据的整理、表示与评价材料构建}
\label{chap:radar-data}

\section{数据工作的目标与已有资产}

本论文将公开数据上的方法验证与雷达领域应用分开。公开数据用于检验执行反馈训练的增量效果；雷达材料用于检查来源约束下的接口适配和问答。领域数据工作的首要问题不是扩大旧图谱的数量，而是明确一条记录描述谁、在什么条件下取何值、依据哪份资料，以及它能用于何种评价。

项目历史上形成了原始文档、网页快照、抽取图谱、检索索引、自动问题和训练轨迹等资产。表\ref{tab:data-legacy}列出已有公开清单及属性剖析中的数量。边数使用 2026 年 9 月 26 日清单中的 21,928 条，不能与更早构建报告中的边数混用；属性剖析绑定相同实体文件的哈希。实体记录还包含制造商、平台和其他对象，10,744 个实体不等于 10,744 个雷达型号。边、属性和过滤后记录也是不同计数口径，不能相加作为独立事实总数。

\begin{table}[htbp]
\centering\small
\caption{历史资产快照的机器统计；数量不表示事实已核验}
\label{tab:data-legacy}
\DataLegacyInventoryTable
\end{table}

旧属性中有 8,470 条包含数值和单位，但这只说明字段存在。剖析对预定的显式条件或限定字段检查得到零条，也不能推断原始资料完全没有条件：条件可能混在属性名称或文本中。由此需要补充结构化条件的保留和来源核验，而不能将所有旧数值直接用作可信标签。

历史自动问题有 train、dev、test 三个标签，数量见表\ref{tab:data-legacy-split}，另有 4,277 条对应训练轨迹。这些材料可用于历史工具诊断，不因带有 test 标签就成为独立人工测试集。清单中的训练与测试锚点交集为零，但从支持字段 \texttt{head/a/b} 提取的 315 个测试显式查询实体中，有 143 个也在训练中；该检查还没有覆盖答案实体和隐式引用。因此原有划分只能支持锚点隔离的描述，不能称严格实体隔离、来源家族隔离或组合泛化。

\begin{table}[htbp]
\centering\small
\caption{旧自动问题的保存标签与数量；不作独立人工评价认定}
\label{tab:data-legacy-split}
\DataLegacySplitTable
\end{table}

\section{四类材料及各自用途}

\subsection{历史库存与旧开发读法}

历史库存用于追溯，不整体升级为已核验知识库。首批审阅从九条旧属性有目的地拆出 12 条候选，涉及六个实体标签和六份二手网页快照，检查主体混淆、事件区别、区间端点、部件标签及未知值等问题。这不是随机样本，不能据其问题比例推断整个旧图谱的错误率。

候选保留原始属性和来源字节，不直接修改旧记录。一个典型风险是检索锚点指向雷达型号，引用页面实际描述整套系统；系统的事件年份不能自动归给其中的雷达。另一些记录涉及把不同部件范围合并、将取消事件写成服役事件，或把来源缺值变成零。这些例子用于规定审阅边界，本章不提供未经独立核验的替代装备参数。

独立上下文的 AI 子代理完成初审后，主代理按开发用途保留两套意见并协调范围。最终开发处置为五条修订或隔离型号断言、五条保留限定来源读法、两条保留读法并继续核对范围。由此形成单独版本的 12 条开发读法及 12 道开发问题。后续中英文改写仍对应相同 12 个目标，每语言 12 题，不是 24 个独立目标；原题、改写和翻译均已暴露，永久排除于最终独立评价。

原人工审阅表仍有 12 条待审，独立事实准入数为零。AI 对来源文字的解读可以支持注明等级的开发实验，但不能代填真人身份或把 \texttt{development\_only} 改成已完成的人工金标。

\subsection{共同接口适配的合成材料}

为学习查询签名而不记忆真实装备答案，构造 100 个虚构知识实例组，每组含一对中英文问题。训练、开发、保留组的分配见表\ref{tab:data-synthetic}。划分先于轨迹生成，实体标识、知识实例和外层措辞模板按组隔离，两种语言留在同组。每个知识实例保留三个虚构实体和 32 条读法，包含竞争字段和条件，不将工具输入缩减成只有目标答案的一条记录。

\begin{table}[htbp]
\centering\small
\caption{合成接口材料的分组；双语视图不是独立事实样本}
\label{tab:data-synthetic}
\DataSyntheticSplitTable
\end{table}

材料覆盖普通属性、条件标量、条件范围、事件、范围和明确缺值等接口情形。真实开发材料中的实体、来源 URI、文本答案和数值答案被排除；共享的属性词、单位、工具名称及条件标签属于有意保留的接口词汇。所有数值都是虚构字段，不保证物理合理性，也不描述真实装备。

目标记录先由声明式实体、属性和条件选择确定，再检查参考程序及成功轨迹与之相同。Agent 每条监督三个成功助手动作，P 监督一次程序输出；工具观察不计入监督目标，不增加失败恢复轨迹。四个 A/C 检查点复用同一 Agent 语料，P 使用相同问题实例。训练效果属于后续实验，本章只说明语料产生方式与数据边界。

这里隔离的是实例和外层措辞，任务家族、属性体系及短工具序列仍然共享。100 组不能写成真实雷达数据库规模，200 个语言视图也不能当成 200 个来源。合成近邻任务上接口可用，与真实资料的事实准确率是不同问题。

\subsection{新制造商来源的探索材料}

另行整理四份制造商 PDF，涉及五个型号、15 个完整页块、62 条结构化来源读法和 24 道中文单意图问题。表\ref{tab:data-sources}列出每份资料的页数和问题数，文档 URL、标题、版本标识及 PDF 哈希保存于公开来源索引。整页抽取在问题窗口之外固定；知识库保留全部型号和干扰读法。每道题在本批只声明一个目标读法，但 62 条读法不只包含 24 条目标记录。

\begin{table}[htbp]
\centering\small
\caption{新来源材料的组成；每份来源内的问题具有相关性}
\label{tab:data-sources}
\DataSourceTable
\end{table}

\begin{table}[htbp]
\centering\small
\caption{新来源读法的值类型；记录数量不等于独立证据数量}
\label{tab:data-values}
\DataValueTypeTable
\end{table}

62 条中包含同一脉宽与 PRF 配置的不同查询方向表示，因此不能视为 62 份互相独立的事实依据。八道题具有预定义的显式查询条件；本批没有未知或事件问答，不能用合成材料的覆盖替代这些真实题型的验证。资料和参考均经 AI 审阅，当前只属于来源约束的探索材料，不是最终独立人工评价集。

\section{记录表示及其实现范围}

\subsection{主体、条件、量值与版本的保留}

将一个可查询读法概括为
\begin{equation}
 r=(e,s,a,k,v,u,c,t,\pi),
\end{equation}
其中 $e$ 是检索锚点，$s$ 是来源所描述的主体，$a$ 是属性，$k$ 是值类型，$v$ 是值或值集合，$u$ 是标准单位，$c$ 是原文条件，$t$ 是事件标签，$\pi$ 是来源定位。分别保存 $e$ 与 $s$，是为了避免在历史锚点归属存在疑问时，将来源描述无条件转换成该型号事实。归属是否可用也独立记录，不因 \texttt{Find} 找到名称就自动确认。

标量、连续范围、离散选项和未知的含义不同。范围保留两个端点；离散选项保留列表含义，不能用最小与最大选项制造连续区间。新来源的选项在知识记录中采用序列化列表字符串，与独立语义断言的列表对应，并未增加集合运算。未知值使用空值和 \texttt{unknown\_scope} 表示资料范围；一个明确缺值记录与查询没有匹配记录的空集合必须区分。事件标签区分服役、初始作战能力、取消和退役等语义，不能仅保留年份。

原始事实模板设有 \texttt{value\_raw}、\texttt{unit\_raw}、\texttt{variant}、条件状态和结构化条件等字段。旧候选及开发版本保留原始字段，新制造商材料通过原文引文、页块和标准单位保存原始表达。但运行接口的必需字段更窄，并不普遍暴露原单位、型号版本或结构化条件数组；不能把规范中的完整字段表写成模型已经可以逐项查询的能力。旧 12 条开发记录的 \texttt{variant} 均为空，型号版本尚未核实，不能因字段存在就认定已经消除版本歧义。

部件或场景标签目前主要保存在来源主体、\texttt{condition\_raw} 及表格上下文中，没有独立的通用部件本体和版本推理模块。文档版本通过资料标题、固定快照和哈希区分，不能等同于装备型号版本已核验。不同来源说法有差异时保留并列记录，不自动选择某个值作为唯一真值。单位只在有依据的资料解读中填写，当前领域执行器不执行通用单位换算。

\subsection{来源定位与可追溯性}

对新 PDF 记录，定位可表示为
\begin{equation}
 \pi=(\mathrm{URI},h_{\mathrm{PDF}},p,h_{\mathrm{text}},[b,e)),
\end{equation}
其中 $p$ 为从 1 开始的 PDF 页号，$h$ 为 SHA-256，$[b,e)$ 为抽取页文本中的 Unicode 码点区间。引用另保留来源 ID、页块 ID、原文件路径和表格上下文。程序检查摘录在对应页文本中恰好出现一次，再计算位置；重复或无法定位的摘录不能凭近似字符串直接通过。

旧网页快照则使用本地文件哈希、JSON pointer、字符区间与已有 chunk 关联。URI 说明来源入口，哈希说明实际保存的字节，二者不可替代。网页更新或抽取文本变更时必须形成新的绑定版本，不能继续使用旧审阅决定。PDF 抽取可能损失列对齐、旋转文字或字体信息，位置一致也不能证明语义解释正确；必要的表格核验仍应回到原 PDF 版面。

新材料分别保存原文件、完整页文本、结构化读法、仅含编号与题面的推理问题，以及单独的参考。公开导出只包含来源元数据、数量、哈希和聚合检查结果，不复制全文、真实题目或逐题答案。数据可追溯性依赖本地保留的原文件；公开哈希本身既不能恢复文档，也不能证明作者具有最终事实核验资格。

\subsection{已实现的只读查询接口}

当前 \texttt{DevelopmentKB} 加载明确标记为非人工、非独立金标的开发读法，检查必要字段、范围端点顺序、未知值范围及引用结构。它提供 \texttt{Find}、\texttt{QueryAttr} 和 \texttt{QueryAttrUnderCondition} 三种查询，后者只对 \texttt{condition\_raw} 或 \texttt{event\_type} 作精确字符串匹配。输出保留不同读法，而不把条件缺失解释成任意条件均成立。

这一接口是来源记录上的可执行查询层，并非将整个历史图谱重建成完整推理图。它没有通用关系遍历、多跳组合、集合运算、区间比较、单位换算或语义条件等价判定。运行层也不会独立鉴定引用主体、核实出版机构身份或消解来源冲突。数值、选项和单位的更广类型检查若存在于回答解析器中，也不能反推知识库执行器拥有对应计算能力。

本章数据组织工作的价值在于为可回放的查询与来源审阅提供清楚边界，尚无证据证明这套表示本身优于其他图谱模式，不把字段整理直接包装为新算法。

\section{构建、审阅与划分流程}

\subsection{语义声明与程序检查的先后关系}

固定来源快照与抽取页后，先声明来源读法及语义参考，再进行引用定位与参考程序验证。另一个 AI 对原 PDF 的主体、量值和表格条件审阅在模型评价之前完成，这不表示所有独立审阅都先于 CPU 检查。新来源的参考与记录先写入独立文件，执行器只检查声明之间是否一致，不能用自身输出制造答案或自动修正标签。\texttt{source\_packet} 同时核对事实 ID、来源、语义字段和查询意图；检查失败即拒绝构建一致的材料包。

\begin{table}[htbp]
\centering\small
\caption{现有数据检查结果及其证据范围}
\label{tab:data-quality}
\resizebox{\linewidth}{!}{\DataQualityTable}
\end{table}

表\ref{tab:data-quality}区分了精确定位、程序自洽和语义审阅。程序成功只能说明已声明的目标可以由该接口获得；引用哈希、JSON 格式和片段位置全部正确，也可能保留主体错配或表格读错。机器检查与语义检查因此分别记录，不能相互代替。

AI 多代理审阅已经用于旧读法、合成材料与新来源包，但各阶段盲性不同。旧候选初审能看到候选备注，主审保有项目上下文；合成审计未见接口适配输出，却已知此前问题；新来源审阅另读 PDF 与表格。这些程序性分工有助于暴露问题，不保证统计独立或不同模型家族的判断。第六章对生成正文另作隐藏路线标签的双 AI 审阅，也仍非人工金标。

原标注规范将人工确认快照读法与独立事实准入分开：前者只确认一份资料的表述，后者还需有实际人工核验人、理由和相应来源支持。软件能检查声明是否完整、引用是否匹配，但无法证明署名者实际完成核验或资料事实绝对正确。当前不能把 AI 审阅通过数量填入人工验收统计。

\subsection{来源家族分组与已暴露材料排除}

后续独立评价应先以来源和型号家族建立关联组。同一原始资料派生的问题、近似改写、翻译以及跨型号共享的来源表格应保持在同组；先分组，再决定开发与评价用途。旧自动问答的完整来源谱系尚未认证，不能仅对题号或字符串去重就宣布来源隔离。原 12 条开发读法及衍生题也不能通过更名、重编号或改写重新进入独立测试。

合成材料已完成按虚构知识实例和措辞模板的分组；新制造商包则与实际调试的旧来源及家族闭包隔离。但 WRM200、METEOR 735C 所属家族和 WR2120 在更早的未使用库存中出现过，基础模型预训练暴露也未知；JRC 型号未在检查的主要旧文件中检出，并非全项目未见的证明。因此四个新来源是本轮探索的相关单元，不能称完全未见家族泛化或最终独立领域测试。

当前 24 道题已经完成评价并暴露，今后只保留为诊断证据。若扩大领域评价，需要在读取新模型成绩之前，预先定义来源多样性、量值形式、事件、条件和部件等覆盖要求，再据资料质量决定具体数量。原计划中的更大规模只是工作量设想，尚未形成数据成果；不能依模型表现选来源、删题或将易答样例补足为所谓独立测试。

\section{与方法章及应用章的衔接}

本章形成四类用途清楚的材料：历史资产用于追溯，旧开发读法用于错误诊断，虚构实例用于共同接口适配，新制造商资料用于有限来源约束评价。第四、五章的公开方法实验不依赖将历史雷达资产认定为真值；第六章分别报告接口使用、记录选择、原文支持、正文与引用，避免把数据检查通过写成回答正确率。

数据层已经实现版本化保存、字段与来源约束、可回放检查及注明身份的 AI 审阅。最终来源覆盖、实际人工核验、完整历史谱系认证和更复杂条件表示仍未完成。后续研究应先补足这些证据边界，不能以继续训练替代数据核验，也不能因本轮材料可运行就宣布整篇论文的领域验证充分。
